\usepackage{array}
\usepackage{listings}
% ================================================================
% Configuración de bloques de código (listings) para el instructivo
% ================================================================
% ===== Configuración básica para listings =====
\usepackage{listings}

\lstdefinestyle{ucaCode}{
  backgroundcolor=\color{backcolour},
  basicstyle=\ttfamily\small,
  keywordstyle=\bfseries\color{codeblue},
  commentstyle=\itshape\color{codegray},
  stringstyle=\color{codered},
  numbers=left,
  numberstyle=\tiny\color{codegray},
  stepnumber=1,
  numbersep=8pt,
  frame=single,
  rulecolor=\color{codegray},
  breaklines=true,
  breakatwhitespace=true,
  tabsize=4,
  showstringspaces=false,
  captionpos=b
}
\lstset{style=ucaCode}


\glssetexpandfield{desc} %Comando para expandir otros comandos en la descripcion de los glosarios
\glssetexpandfield{name} %Comando para expandir otros comandos en el nombre de los glosarios
%==========================CONTADORES==============================%
%Se crea nuevo contador que se utiliza para agregar las letras correspondientes a los anexos
\newcounter{anexosletra}

%Se establece 1 como el valor inicial del contador
\setcounter{anexosletra}{1}

% Ancho fijo para la etiqueta "ANEXO X."
\newlength{\anexolabelwidth}
\settowidth{\anexolabelwidth}{ANEXO MM.}

%Se crea el contador que se utiliza para agregar el numero correlativo de las secciones dentro de un anexo
\newcounter{anexosseccion}

%Se crea el contador que se utiliza para agregar el numero correlativo de una subseccion dentro de un anexo
\newcounter{anexossubseccion}

%Se crea el contador que se utiliza para llevar la cuenta de cuantos autores han sido ingresados en la primera portada, esto sirve para saber si la mayoria son hombres o mujeres y tambien para establecer un maximo de cuatro autores
\newcounter{contadorautor}

%Se establece el valor inicial 0 para el contador de autores
\setcounter{contadorautor}{0}

%Se crea el contador que se utiliza para llevar la cuenta de cuantos autores son hombres
\newcounter{contadorhombrescarrerauno}

%Se establece el valor inicial 0 para el contador de hombres
\setcounter{contadorhombrescarrerauno}{0}

%Se crea el contador que se utiliza para llevar la cuenta de cuantos autores son mujeres
\newcounter{contadormujerescarrerauno}

%Se establece el valor inicial 0 para el contador de mujeres
\setcounter{contadormujerescarrerauno}{0}

%Se crea el contador que se utiliza para llevar la cuenta de cuantos autores son hombres
\newcounter{contadorhombrescarrerados}

%Se establece el valor inicial 0 para el contador de hombres
\setcounter{contadorhombrescarrerados}{0}

%Se crea el contador que se utiliza para llevar la cuenta de cuantos autores son mujeres
\newcounter{contadormujerescarrerados}

%Se establece el valor inicial 0 para el contador de mujeres
\setcounter{contadormujerescarrerados}{0}

%Se crea el contador que se utiliza para llevar la cuenta de cuantos autores son hombres
\newcounter{contadorhombrescarreratres}

%Se establece el valor inicial 0 para el contador de hombres
\setcounter{contadorhombrescarreratres}{0}

%Se crea el contador que se utiliza para llevar la cuenta de cuantos autores son mujeres
\newcounter{contadormujerescarreratres}

%Se establece el valor inicial 0 para el contador de mujeres
\setcounter{contadormujerescarreratres}{0}

%Se crea el contador que se utiliza para llevar la cuenta de cuantos autores son mujeres
\newcounter{resultadosexosautorescarrerauno}
\setcounter{resultadosexosautorescarrerauno}{0}

\newcounter{resultadosexosautorescarrerados}
\setcounter{resultadosexosautorescarrerados}{0}

\newcounter{resultadosexosautorescarreratres}
\setcounter{resultadosexosautorescarreratres}{0}

\newcounter{figuras}
\setcounter{figuras}{0}

\newcounter{nomenclatura}
\setcounter{nomenclatura}{0}
%===========================COMANDOS=============================%

%Este comando verifica que el parametro que le pasaron este vacio o no
\newcommand{\verificarvacio}[3]{
    \ifdefined #1
        \ifthenelse{\equal{#1}{}}{%
        %Si esta vacio muestra un mensaje de advertencia en mayuscula y de color rojo
        \textcolor{red}{\MakeUppercase{#2}}%
        }{
        %Si no esta vacio entonces muestra el texto del parametro en mayuscula
        \MakeUppercase{#1}
        }
    \else
        \textcolor{red}{\MakeUppercase{#3}}
    \fi
}

% Estos comandos guardan los nombres de los integrantes del grupo, el comando autorA corresponde al primer autor, autorB corresponse al segundo autor, autorC corresponde al tercer autor y autorD corresponde al cuarto autor
\newcommand{\autorAnombre}{}
\newcommand{\autorAcarrera}{}
\newcommand{\autorBnombre}{}
\newcommand{\autorBcarrera}{}
\newcommand{\autorCnombre}{}
\newcommand{\autorCcarrera}{}
\newcommand{\autorDnombre}{}
\newcommand{\autorDcarrera}{}

%El comando autor se encarga de agregar los nombres de los autores en los comandos anteriores
\newcommand{\autor}[3]{
    %Se verifica cuantas veces se a utilizado el comando autor, si se ha hecho menos de cuatro veces entonces todavia se pueden agregar mas autores, si se ha hecho igual o mas de cuatro veces entonces ya no se puede agregar mas autores
    \ifthenelse{\value{contadorautor} < 4}{
        %El contador "contadorautor" sirve para llevar la cuenta de cuantos autores se han ingresado, el comando \ifcase sirve para verificar que valor tiene el contador en ese momento y en base a eso hacer una acción u otra.
        \ifcase\value{contadorautor}%
            %Si el contador es 0 no se ha ingresado ningun autor todavia, por lo tanto se agrega el nombre al comando autorA
            \renewcommand{\autorAnombre}{#1}
            \renewcommand{\autorAcarrera}{#3}
        \or%
            %Si el contador es 1 se ha ingresado ya un autor, por lo tanto se agrega el nombre al comando autorB
            \renewcommand{\autorBnombre}{#1}
            \renewcommand{\autorBcarrera}{#3}
        \or%
            %Si el contador es 2 se han ingresado dos autores, por lo tanto se agrega el nombre al comando autorC
            \renewcommand{\autorCnombre}{#1}
            \renewcommand{\autorCcarrera}{#3}
        \or%
            %Si el contador es 3 se han ingresado tres autores, por lo tanto se agrega el nombre al comando autorD
            \renewcommand{\autorDnombre}{#1}
            \renewcommand{\autorDcarrera}{#3}
        \fi

        \ifthenelse{\equal{#3}{carreraUno}}{
        %Se valida que el sexo ingresado para cada un autor sea hombre o mujer, si es mujer el contador "contadormujeres" incrementa en 1, si en hombre el contador "contadorhombres incrementa en 1
        \ifthenelse{\equal{#2}{M}}{
            \stepcounter{contadormujerescarrerauno}
        }{
            \stepcounter{contadorhombrescarrerauno}
        }
        }{
            \ifthenelse{\equal{#3}{carreraDos}}{
            %Se valida que el sexo ingresado para cada un autor sea hombre o mujer, si es mujer el contador "contadormujeres" incrementa en 1, si en hombre el contador "contadorhombres incrementa en 1
            \ifthenelse{\equal{#2}{M}}{
                \stepcounter{contadormujerescarrerados}
            }{
                \stepcounter{contadorhombrescarrerados}
            }
            }{
                \ifthenelse{\equal{#3}{carreraTres}}{
                %Se valida que el sexo ingresado para cada un autor sea hombre o mujer, si es mujer el contador "contadormujeres" incrementa en 1, si en hombre el contador "contadorhombres incrementa en 1
                \ifthenelse{\equal{#2}{M}}{
                    \stepcounter{contadormujerescarreratres}
                }{
                    \stepcounter{contadorhombrescarreratres}
                }
                }{}
            }
        }

        %Se incrementa en 1 el contador "contadorautor" al final del comando
        \stepcounter{contadorautor}
    }{
    
    }
}

%Este comando sirve para verificar el numero de mujeres y hombres en el grupo de trabajo ya que es necesario para escribir el grado a optar en la primera portada del trabajo
\newcommand{\verificarsexoautores}{
    \ifthenelse{\equal{\value{contadormujerescarrerauno}}{0}}{
        %Si no hay mujeres en el grupo entonces el valor del contador resultadosexosautores es 1
        \setcounter{resultadosexosautorescarrerauno}{1}
    }{
        \ifthenelse{\equal{\value{contadorhombrescarrerauno}}{0}}{
            %Si no hay hombres en el grupo entonces el valor del contador resultadosexosautores es 2
            \setcounter{resultadosexosautorescarrerauno}{2}
        }{
            \ifthenelse{\value{contadorhombrescarrerauno} > \value{contadormujerescarrerauno} }{
                %Si hay mas hombres que mujeres entonces el valor de resultadosexosautores es 3
                \setcounter{resultadosexosautorescarrerauno}{3}
            }{
                \ifthenelse{\value{contadormujerescarrerauno} > \value{contadorhombrescarrerauno}}{
                    %Si hay mas mujeres que hombres entonces el valor de resultadosexosautores es 4
                    \setcounter{resultadosexosautorescarrerauno}{4}
                }{
                    \ifthenelse{\equal{\value{contadorhombrescarrerauno}}{\value{contadormujerescarrerauno}}}{
                        %Si hay igual numero de mujeres que de hombres el valor del contadore resultadosexosautores es 5
                        \setcounter{resultadosexosautorescarrerauno}{5}
                    }{
                    
                    }
                }
            }
        }
    }
    \ifthenelse{\equal{\value{contadormujerescarrerados}}{0}}{
        %Si no hay mujeres en el grupo entonces el valor del contador resultadosexosautores es 1
        \setcounter{resultadosexosautorescarrerados}{1}
    }{
        \ifthenelse{\equal{\value{contadorhombrescarrerados}}{0}}{
            %Si no hay hombres en el grupo entonces el valor del contador resultadosexosautores es 2
            \setcounter{resultadosexosautorescarrerados}{2}
        }{
            \ifthenelse{\value{contadorhombrescarrerados} > \value{contadormujerescarrerados} }{
                %Si hay mas hombres que mujeres entonces el valor de resultadosexosautores es 3
                \setcounter{resultadosexosautorescarrerados}{3}
            }{
                \ifthenelse{\value{contadormujerescarrerados} > \value{contadorhombrescarrerados}}{
                    %Si hay mas mujeres que hombres entonces el valor de resultadosexosautores es 4
                    \setcounter{resultadosexosautorescarrerados}{4}
                }{
                    \ifthenelse{\equal{\value{contadorhombrescarrerados}}{\value{contadormujerescarrerados}}}{
                        %Si hay igual numero de mujeres que de hombres el valor del contadore resultadosexosautores es 5
                        \setcounter{resultadosexosautorescarrerados}{5}
                    }{
                    
                    }
                }
            }
        }
    }
    \ifthenelse{\equal{\value{contadormujerescarreratres}}{0}}{
        %Si no hay mujeres en el grupo entonces el valor del contador resultadosexosautores es 1
        \setcounter{resultadosexosautorescarreratres}{1}
    }{
        \ifthenelse{\equal{\value{contadorhombrescarreratres}}{0}}{
            %Si no hay hombres en el grupo entonces el valor del contador resultadosexosautores es 2
            \setcounter{resultadosexosautorescarreratres}{2}
        }{
            \ifthenelse{\value{contadorhombrescarreratres} > \value{contadormujerescarreratres} }{
                %Si hay mas hombres que mujeres entonces el valor de resultadosexosautores es 3
                \setcounter{resultadosexosautorescarreratres}{3}
            }{
                \ifthenelse{\value{contadormujerescarreratres} > \value{contadorhombrescarreratres}}{
                    %Si hay mas mujeres que hombres entonces el valor de resultadosexosautores es 4
                    \setcounter{resultadosexosautorescarreratres}{4}
                }{
                    \ifthenelse{\equal{\value{contadorhombrescarreratres}}{\value{contadormujerescarreratres}}}{
                        %Si hay igual numero de mujeres que de hombres el valor del contadore resultadosexosautores es 5
                        \setcounter{resultadosexosautorescarreratres}{5}
                    }{
                    
                    }
                }
            }
        }
    }
}

%Valor booleano para saber si ya se ingreso alguna entrada a las siglas o no
\newcommand{\ningunasiglaingresada}{true} 

%Este comando sirve para agregar nuevas siglas a la seccion de siglas
\newcommand{\itemsiglas}[2]{%
    %Valida que se agreguen entradas a las siglas solo si los dos parametros del comando \itemsiglas han sido llenados, si no han sido llenados no se agreaga ninguna sigla nueva
    \ifthenelse{\equal{#1}{} \or \equal{#2}{}}{ 
    
    }{
    %Si los dos parametros fueron agregados entonces se ingresa la nueva sigla y \ningunasiglaingresada cambia a falso
        \renewcommand{\ningunasiglaingresada}{false}
        \newglossaryentry{#1}{name={#1}, type=siglas, description={#2}}%
    }
}

%Valor booleano para saber si ya se ingreso alguna entrada a las abreviaturas o no
\newcommand{\ningunaabreviaturaingresada}{true} 

%Este comando sirve para agregar nuevas abreviaturas a la sección de abreviaturas
\newcommand{\itemabreviatura}[2]{
    %Valida que se agreguen entradas a las abreviaturas solo silos dos parametros del comando \itemabreviatura han sido llenados, si no han sido llenados no se agreaga ninguna abreviatura nueva
    \ifthenelse{\equal{#1}{} \or \equal{#2}{}}{  
    }{
    %Si los dos parametros fueron agregados entonces se ingresa la nueva abreviatura y \ningunaabreviaturaingresada cambia a falso
        \renewcommand{\ningunaabreviaturaingresada}{false}
        \newglossaryentry{#1}{name={#1}, type=abreviaturas, description={#2}}%
    }
    
}

%Valor booleano para saber si ya se ingreso alguna entrada a la nomenclatura o no
\newcommand{\ningunanomenclaturaingresada}{true}

%Este comando sirve para agregar un nuevo elemento a la sección de nomenclatura
\newcommand{\itemnomenclatura}[2]{%
    %Valida que se agreguen entradas a la nomenclatura solo silos dos parametros del comando \itemnomenclatura han sido llenados, si no han sido llenados no se agreaga ninguna nomenclatura nueva
    \ifthenelse{\equal{#1}{} \or \equal{#2}{}}{       
    }{
    \stepcounter{nomenclatura}
    %Si los dos parametros fueron agregados entonces se ingresa la nueva nomenclatura y \ningunanomenclaturaingresada cambia a falso
        \renewcommand{\ningunanomenclaturaingresada}{false}
        \newglossaryentry{#2}{name={#1}, type=nomenclatura, description={#2}}%
    }
}

%Valor booleano para saber si ya se ingreso alguna entrada al glosario o no
\newcommand{\glosariovacio}{true} 

%Este comando sirve para agregar nuevas palabras al glosario
\newcommand{\itemglosario}[2]{%
    %Valida que se agreguen entradas al glosario solo si los dos parametros del comando \glosario han sido llenados, si no han sido llenados no se agreaga ninguna entrada al glosario nueva
    \ifthenelse{\equal{#1}{} \or \equal{#2}{}}{ 
        
    }{
    %Si los dos parametros fueron agregados entonces se ingresa la nueva
    %entrada del glosario y \glosariovacio cambia a falso
        \renewcommand{\glosariovacio}{false}
        \newglossaryentry{#1}{name={#1}, type=glosario, description={#2}}%
    }
}


%Este comando esta creado para hacer uno o dos saltos de pagina dependiendo de si la seccion termina en pagina impar o par
\newcommand{\paroimpar}[1]{%
    \clearpage
    \ifodd\value{page}% ya estamos en impar, no hacer nada
    \else
        \hbox{}\thispagestyle{empty}% página en blanco real (sin numeración)
        \clearpage
    \fi
}


%Comando utilizado para crear los capitulos
\newcommand{\capitulo}[2]{
    \begin{capituloentorno}
    {
        #1
    }{
        #2
    }
    \end{capituloentorno}
}

%Valor booleano para saber si ya se agrego una imagen al documento o no
\newcommand{\ningunaimageningresada}{true}

% Se crea nuevo comando que establece el formato de las figuras
\newcommand{\imagen}[6]{
    \renewcommand{\ningunaimageningresada}{false}
    %Se valida si existe alguna imagen con el nombre proporcionado
    \IfFileExists{img/#1}{
        %Se cambia el valor del comando ningunaimageningresada a false
        \renewcommand{\ningunaimageningresada}{false}
        %\stepcounter{figuras}
        %Si el cuarto parametro del comando es V significa que se quiere mostrar la imagen en formato vertical.
        \ifthenelse{\equal{#6}{V}}{
            %Se valida que el tercer parametro no este vacio, el tercer parametro corresponde a la escala que tendra la figura, si el parametro esta vacío entonces se coloca 1 por defecto
            \ifthenelse{\equal{#5}{}}{
                \begin{figure}[H]
                \vspace{2\baselineskip}
                \centering
                \includegraphics[scale=1]{img/#1}
                \ifthenelse{\equal{#3}{true}}{
                    \caption[#2]{#2. Adaptado de: [#4]}
                }{
                    \caption[#2]{#2. Fuente: [#4]}
                }
                \label{fig:#1}
                \end{figure}
            }{
            %Si no esta vacio entonces se coloca el valor que se escribio en el tercer parametro
                \begin{figure}[H]
                \vspace{2\baselineskip}
                \centering
                \includegraphics[scale=#5]{img/#1}
                \ifthenelse{\equal{#3}{true}}{
                    \caption[#2]{#2. Adaptado de: [#4]}
                }{
                    \caption[#2]{#2. Fuente: [#4]}
                }
                \label{fig:#1}
                \end{figure}
            }    
        }{
        %Si el cuarto parametro es H significa que la figura se mostrara de manera horizontal
        \ifthenelse{\equal{#6}{H}}{
            \begin{sidewaysfigure}
                \includegraphics[width=\linewidth]{img/#1}
                \ifthenelse{\equal{#3}{true}}{
                    \caption[#2]{#2. Adaptado de: [#4]}
                }{
                    \caption[#2]{#2. Fuente: [#4]}
                }
                \label{fig:#1}
            \end{sidewaysfigure}
        }{
        \textcolor{red}{\MakeUppercase{ERROR: el parámetro para la orientación de la imagen solo puede ser 'H' para horizontal y 'V' para vertical}}
        }
        }
    }{
        \textcolor{red}{\MakeUppercase{ERROR: la imagen con nombre #1 no existe}}
    }
}

%Este comando sirve para escribir la carrera que se esta cursando de la manera correcta en la primera porta validando la cantidad de mujeres y hombres en el grupo
\newcommand{\determinarcarrera}[3]{
    %Si el segundo parametro del comando es 1 significa que se inserata el nombre de la carrera en la segunda portada en la parte del director de carrera
    \ifthenelse{\equal{#2}{1}}{%
    {%
      \centering
      \ifdefstring{#1}{informatica}{INGENIERÍA INFORMÁTICA\par}{%
      \ifdefstring{#1}{civil}{INGENIERÍA CIVIL\par}{%
      \ifdefstring{#1}{electrica}{INGENIERÍA ELÉCTRICA\par}{%
      \ifdefstring{#1}{energetica}{INGENIERÍA ENERGÉTICA\par}{%
      \ifdefstring{#1}{alimentos}{INGENIERÍA DE ALIMENTOS\par}{%
      \ifdefstring{#1}{quimica}{INGENIERÍA QUÍMICA\par}{%
      \ifdefstring{#1}{mecanica}{INGENIERÍA MECÁNICA\par}{%
      \ifdefstring{#1}{industrial}{INGENIERÍA INDUSTRIAL\par}{%
      \ifdefstring{#1}{arquitectura}{ARQUITECTURA\par}{%
      \ifdefstring{#1}{diseño}{LICENCIATURA EN DISEÑO\par}{%
      {\textcolor{red}{ERROR: NOMBRE DE CARRERA INVÁLIDO}\par}%
      }}}}}}}}}}}%
    }%
    {%
    %Si el segundo parametro no es 1 significa que se esta insertando el nombre de la carrera en la parte donde se menciona el grado que se quiere obtener en la primera portada

    %Si el valor del contador resultadosexosautores es 1 entonces solamente hay hombres en el grupo y se pone de la siguiente manera
    \ifthenelse{\equal{#3}{1}}{
        {%
        \centering
        \ifdefstring{#1}{informatica}{INGENIERO INFORMÁTICO\par}{%
        \ifdefstring{#1}{civil}{INGENIERO CIVIL\par}{%
        \ifdefstring{#1}{electrica}{INGENIERO ELÉCTRICO\par}{%
        \ifdefstring{#1}{energetica}{INGENIERO ENERGÉTICO\par}{%
        \ifdefstring{#1}{alimentos}{INGENIERO DE ALIMENTOS\par}{%
        \ifdefstring{#1}{quimica}{INGENIERO QUÍMICO\par}{%
        \ifdefstring{#1}{mecanica}{INGENIERO MECÁNICO\par}{%
        \ifdefstring{#1}{industrial}{INGENIERO INDUSTRIAL\par}{%
        \ifdefstring{#1}{arquitectura}{ARQUITECTO\par}{%
        \ifdefstring{#1}{diseño}{LICENCIADO EN DISEÑO\par}{%
        {\textcolor{red}{ERROR: NO SE HA INGRESADO UN NOMBRE DE CARRERA VÁLIDO}\par}%
        }}}}}}}}}}}%
    }%
    {
            %Si el valor del contador "resultadosexosautores" es 2 entonces solamente hay mujeres en el grupo y se pone de la siguiente manera
            \ifthenelse{\equal{#3}{2}}{
            {%
                \centering
                \ifdefstring{#1}{informatica}{INGENIERA INFORMÁTICA\par}{%
                \ifdefstring{#1}{civil}{INGENIERA CIVIL\par}{%
                \ifdefstring{#1}{electrica}{INGENIERA ELÉCTRICA\par}{%
                \ifdefstring{#1}{energetica}{INGENIERA ENERGÉTICA\par}{%
                \ifdefstring{#1}{alimentos}{INGENIERA DE ALIMENTOS\par}{%
                \ifdefstring{#1}{quimica}{INGENIERA QUÍMICA\par}{%
                \ifdefstring{#1}{mecanica}{INGENIERA MECÁNICA\par}{%
                \ifdefstring{#1}{industrial}{INGENIERA INDUSTRIAL\par}{%
                \ifdefstring{#1}{arquitectura}{ARQUITECTA\par}{%
                \ifdefstring{#1}{diseño}{LICENCIADA EN DISEÑO\par}{%
                {\textcolor{red}{ERROR: NO SE HA INGRESADO UN NOMBRE DE CARRERA VÁLIDO}\par}%
                }}}}}}}}}}}%
        }{
            %Si hay mas hombres que mujeres en el grupo entonces el valor del contador "resultadosexosautores" es 3 y se escribe de la siguiente manera 
            \ifthenelse{\equal{#3}{3} \or \equal{#3}{5}}{
            {%
                \centering
                \ifdefstring{#1}{informatica}{INGENIERO(A) INFORMÁTICO(A)\par}{%
                \ifdefstring{#1}{civil}{INGENIERO(A) CIVIL\par}{%
                \ifdefstring{#1}{electrica}{INGENIERO(A) ELÉCTRICO(A)\par}{%
                \ifdefstring{#1}{energetica}{INGENIERO(A) ENERGÉTICO(A)\par}{%
                \ifdefstring{#1}{alimentos}{INGENIERO(A) DE ALIMENTOS\par}{%
                \ifdefstring{#1}{quimica}{INGENIERO(A) QUÍMICO(A)\par}{%
                \ifdefstring{#1}{mecanica}{INGENIERO(A) MECÁNICO(A)\par}{%
                \ifdefstring{#1}{industrial}{INGENIERO(A) INDUSTRIAL\par}{%
                \ifdefstring{#1}{arquitectura}{ARQUITECTO(A)\par}{%
                \ifdefstring{#1}{diseño}{LICENCIADO(A) EN DISEÑO\par}{%
                {\textcolor{red}{ERROR: NO SE HA INGRESADO UN NOMBRE DE CARRERA VÁLIDO}\par}%
                }}}}}}}}}}}%
            }{
                %Si hay mas mujeres que hombres entonces el valor del contador "resultadosexosautores" es 4 y se escribe de la siguiente forma
                \ifthenelse{\equal{#3}{4}}{%
                    {%
                    \centering
                    \ifdefstring{#1}{informatica}{INGENIERA(O) INFORMÁTICA(O)\par}{%
                    \ifdefstring{#1}{civil}{INGENIERA(O) CIVIL\par}{%
                    \ifdefstring{#1}{electrica}{INGENIERA(O) ELÉCTRICA(O)\par}{%
                    \ifdefstring{#1}{energetica}{INGENIERA(O) ENERGÉTICA(O)\par}{%
                    \ifdefstring{#1}{alimentos}{INGENIERA(O) DE ALIMENTOS\par}{%
                    \ifdefstring{#1}{quimica}{INGENIERA(O) QUÍMICA(O)\par}{%
                    \ifdefstring{#1}{mecanica}{INGENIERA(O) MECÁNICA(O)\par}{%
                    \ifdefstring{#1}{industrial}{INGENIERA(O) INDUSTRIAL\par}{%
                    \ifdefstring{#1}{arquitectura}{ARQUITECTA(O)\par}{%
                    \ifdefstring{#1}{diseño}{LICENCIADA(O) EN DISEÑO\par}{%
                    {\textcolor{red}{ERROR: NO SE HA INGRESADO UN NOMBRE DE CARRERA VÁLIDO}\par}%
                    }}}}}}}}}}}%
                }{}
                }
            }
        }
    } 
}

%Valor booleano para validar si ninguna tabla ha sido ingresada
\newcommand{\ningunatablaingresada}{true}

%Este comando sirve para insertar una tabla en el documento utlizando el formato requerido
\newcommand{\tabla}[2]{
    \renewcommand{\ningunatablaingresada}{false}

    % V = vertical (normal)
    \ifthenelse{\equal{#1}{V}}{
        #2
    }{
    % H = horizontal (antes: landscape). En Arquitectura: usar ancho tabloide completo
    \ifthenelse{\equal{#1}{H}}{
        \ifarqformato
            \clearpage
            \onecolumn
            #2
            \clearpage
            \twocolumn
            \setlength{\columnsep}{1cm}%
        \else
            \begin{landscape}
                #2
            \end{landscape}
        \fi
    }{
        \textcolor{red}{\MakeUppercase{ERROR: el parametro para la orientacion de la tabla solo puede ser 'H' para horizontal y 'V' para vertical}}
    }
    }
}


%Este comando se encarga de crear la portada de un anexo
%Este comando se encarga de crear la portada de un anexo
\newcommand{\portadanexo}[1]{%
  % Terminar bien lo anterior
  \clearpage

  % Arquitectura: pasar a una columna ANTES de asegurar impar
  \ifarqformato
    \onecolumn
  \fi

  % Asegurar que la portada del anexo inicie en página impar
  \cleardoublepage

  \thispagestyle{empty}

  % ===== Bloque centrado verticalmente =====
\vspace*{\fill}

\noindent\makebox[\textwidth][c]{%
  \parbox{%
    \ifarqformato 0.60\textwidth \else 0.95\textwidth \fi
  }{%
    \centering
    {\fontsize{20}{24}\selectfont
      ANEXO \MakeUppercase{\alph{anexosletra}}\par
    }
    %\vspace*{0.7\baselineskip}
    {\fontsize{16}{24}\selectfont
      \MakeUppercase{#1}\par
    }
  }%
}

\vspace*{\fill}


  % =========================================

  % Pasar a la siguiente página
  \newpage

  % Página en blanco detrás de la portada (sin numeración visible)
  \thispagestyle{empty}\null

  % Arquitectura: volver a doble columna para el cuerpo
  \ifarqformato
    \twocolumn
    \setlength{\columnsep}{1cm}%
  \else
    \newpage
  \fi
}



%Este comando crea el cuerpo de un anexo
\newenvironment{cuerpoanexo}[1]{%
    % Para Arquitectura, el cuerpo del anexo vuelve a doble columna
    \ifarqformato
      \twocolumn
      \setlength{\columnsep}{1cm}%
    \fi
    \setlength{\parskip}{\baselineskip}
    \setcounter{page}{1}
    \renewcommand{\thepage}{\MakeUppercase{\alph{anexosletra}}-\arabic{page}}
    #1
    \clearpage
}{%
   
}


\newcommand{\anexo}[2]{
    %\chapter*{\phantom{Anexo \theanexosletra}}
    \hypertarget{anexo:\Alph{anexosletra}}{}


    \setcounter{numeroimagenesanexos}{1}
    \setcounter{anexosseccion}{0}
    
    \addtocontents{toc}{%
      \protect\hyperlink{anexo:\Alph{anexosletra}}{%
        \makebox[\anexolabelwidth][l]{ANEXO \MakeUppercase{\alph{anexosletra}}.}%
        \parbox[t]{380pt}{\strut#1\strut}%
      }\\
    }

    \portadanexo{#1}
    \begin{cuerpoanexo}{
        #2
    }
    \end{cuerpoanexo}

    \stepcounter{anexosletra}
}



%Este comando sirve para agregar imagenes en los anexos utilizando el formato adecuado
\newcommand{\imagenanexo}[4]{
    \IfFileExists{img/#1}{
        \ifthenelse{\equal{#4}{V}}{
            \ifthenelse{\equal{#3}{}}{
                \begin{figure}[H]
                \vspace{2\baselineskip}
                \centering
                \includegraphics[scale=1]{img/#1}
                \caption*{Figura \Alph{anexosletra}.\arabic{numeroimagenesanexos} #2}
                \label{fig:#1}
                \end{figure}
            }{
                \begin{figure}[H]
                \vspace{2\baselineskip}
                \centering
                \includegraphics[scale=#3]{img/#1}
                \caption*{Figura \Alph{anexosletra}.\arabic{numeroimagenesanexos} #2}
                \label{fig:#1}
                \end{figure}
            }    
        }{
        \ifthenelse{\equal{#4}{H}}{
            \begin{sidewaysfigure}
                \centering
                \includegraphics[width=\linewidth]{img/#1}
                \caption*{Figura \Alph{anexosletra}.\arabic{numeroimagenesanexos} #2}
                \label{fig:#1}
            \end{sidewaysfigure}
        }{
        \textcolor{red}{\MakeUppercase{ERROR: el parametro para la orientacion de la imagen solo puede ser 'H' para horizontal y 'V' para vertical}}
        }
        }
        \stepcounter{numeroimagenesanexos}
    }{
        \textcolor{red}{\MakeUppercase{ERROR: la imagen con nombre #1 no existe}}
    }
}

%Este comando sirve para hacer referencia a una figura utilizando el formato correcto
\newcommand{\reffig}[1]{
    Figura \ref{fig:#1}%
}

%Este comando sirve para hace referencia a una tabla utilizando el formato correcto
\newcommand{\reftab}[1]{
    Tabla \ref{#1}%
}

%Este comando permite agregar secciones en los anexos con el formato correcto
\newcommand{\sectionanexo}[1]{
    \stepcounter{anexosseccion}%\\
    \textbf{\Alph{anexosletra}.\theanexosseccion \, #1}%\\
    %\\
    \setcounter{anexossubseccion}{0}
}

%Este comando sirve para escribir una subseccion en los aexos utilizando el formato adecuado
\newcommand{\subsectionanexo}[1]{
    \stepcounter{anexossubseccion}
    \textbf{\Alph{anexosletra}.\theanexosseccion.\theanexossubseccion \, #1}%\\
    %\\
}

%Este comando sirve para escribir la fuente de las tablas de manera apropiada
\newcommand{\fuentetabla}[1]{
    \caption*{Fuente: [#1]}
}

%Este comando sirve para escribir el epigrafe de una tabla que continua en la pagina siguiente y por lo tanto se debe colocar la palabra continución
\newcommand{\captioncontinuacion}[1]{
    \caption*{Tabla \thetable \, #1 (continuación)}
}

%Este comando sirve para mostrar codgio
\newcommand{\codigo}[1]{%
  \begin{quote}%
    \ttfamily%
    \codigoprocess#1\\\relax%
  \end{quote}%
}

%Este comando sirve para mostrar codgio
\def\codigoprocess#1\\#2\relax{%
  \detokenize{#1}%
  \def\temp{#2}%
  \ifx\temp\empty%
  \else%
    \newline%
    \codigoprocess#2\relax%
  \fi%
}

%\usepackage{showframe} % Paquete para mostrar márgenes del documento

\newcolumntype{C}[1]{>{\centering\arraybackslash}p{#1}}
% ================================================================
%            Sección editable para primera portada
% ================================================================
% =============== NO MODIFICAR LOS COMANDOS ======================
% ================================================================

\newcommand{\titulo}{Sistema de Generación Aumentada por Recuperación (RAG) para Diagnóstico y Mantenimiento de Aires Acondicionados tipo Split} %Título del trabajo

% Carrera - [comando correspondiente]:
% 1. Ingeniería Informática - informatica
% 2. Ingeniería Civil - civil
% 3. Ingeniería Eléctrica - electrica
% 4. Ingeniería Energética - energetica
% 5. Ingeniería de Alimentos - alimentos
% 6. Ingeniería Química -quimica
% 7. Ingeniería Mecanica - mecanica
% 8. Ingeniería Industrial - industrial
% 9. Arquitectura - arquitectura
% 10. Licenciatura en Diseño - diseño

\newcommand{\carreraUno}{informatica}
%\newcommand{\carreraDos}{energetica}
%\newcommand{\carreraTres}{industrial}

\autor{Rodrigo Andrés Mena Caballero}{M}{carreraUno} % Autor 1
\autor{Duglas Fernando Pineda Fuentes}{M}{carreraUno} % Autor 2 (opcional)
\autor{Mario Antonio Martinez Villatoro}{M}{carreraUno} % Autor 3 (opcional)
\autor{}{}{} % Autor 4 (opcional)

\newcommand{\mesdegraduacion}{Abril 2026} % Mes y año de graduación

%================================================================
%            Sección editable para segunda portada
% ===============================================================
% ================ NO MODIFICAR LOS COMANDOS ====================
% ===============================================================

\newcommand{\rector}{Mario Ernesto Cornejo Mena, S.J.}% Nombre del rector
\newcommand{\sexoRector}{masculino} % Sexo del secretario o secretaria general
\newcommand{\secretarioGeneral}{Lidia Gabriela Bolaños Teodoro} % Nombre del secretario o secretaria general
\newcommand{\sexoSecretario}{femenino} % Sexo del secretario o secretaria general
\newcommand{\decano}{Carlos Ernesto Rivas Cerna} % Nombre del decano de la facultad de ingeniería y arquitectura
\newcommand{\sexoDecano}{masculino} % Sexo del secretario o secretaria general
\newcommand{\directorDeCarrera}{José Enmanuel Amaya Araújo} % Director de la carrera
\newcommand{\carreraDirectorDeCarrera}{informatica} % Carrera del director de carrera
\newcommand{\sexoDirectorDeCarrera}{masculino} % Sexo del director de carrera
% Segundo director de carrera (opcional, para segunda carrera)
\newcommand{\directorDeCarreraDos}{} % Nombre del segundo director de carrera (dejar vacío si no aplica)
\newcommand{\carreraDirectorDeCarreraDos}{} % Carrera del segundo director
\newcommand{\sexoDirectorDeCarreraDos}{} % Sexo del segundo director de carrera
\newcommand{\directorDeTrabajo}{James Edward Humberstone Morales} % Director del trabajo de graduación
\newcommand{\sexoDirectorDeTrabajo}{masculino} % Sexo del secretario o secretaria general
\newcommand{\lector}{nombre del lector} % Nombre del lector o lectora
\newcommand{\sexoLector}{masculino} % Sexo del lector o lectora

%=================================================================
%            Sección editable para agradecimientos
% ================================================================
% ================ NO MODIFICAR LOS COMANDOS =====================
% ================================================================

\newcommand{\agradecimiento}{
Lorem ipsum dolor sit amet, consectetur adipiscing elit. Mauris eget arcu id justo mollis sagittis. Integer commodo vitae libero eget sagittis. Etiam id ante dignissim, sollicitudin odio id, rhoncus magna. Vestibulum ante ipsum primis in faucibus orci luctus et ultrices posuere cubilia curae; Praesent at tortor eros. Aenean in euismod mauris. Suspendisse pharetra aliquet ex, sed tempus augue aliquam nec. Donec tempor molestie mi, quis faucibus sem. Cras ornare bibendum nulla a laoreet. Vestibulum risus massa, feugiat at eros ut, rhoncus aliquam turpis.
} % Agradecimientos

%=================================================================
%               Sección editable para dedicatorias
% ================================================================
% ================= NO MODIFICAR LOS COMANDOS ====================
% ================================================================

\newcommand{\dedicatoriaautorA}{}% Dedicatoria del autor A, ingresar solo si se ingreso "true" en el comando "porcadaestudiante".

\newcommand{\dedicatoriaautorB}{}%}% Dedicatoria del autor B, ingresar solo si se ingreso "false" en el comando "porcadaestudiante".

\newcommand{\dedicatoriaautorC}{}%}% Dedicatoria del autor C, ingresar solo si se ingreso "false" en el comando "porcadaestudiante".

\newcommand{\dedicatoriaautorD}{}%}% Dedicatoria del autor D, ingresar solo si se ingreso "false" en el comando "porcadaestudiante".

%=================================================================
%                   Sección editable para resumen
% ================================================================
% ================== NO MODIFICAR LOS COMANDOS ===================
% ================================================================

\newcommand{\resumen}{
Lorem ipsum dolor sit amet, consectetur adipiscing elit. Mauris eget arcu id justo mollis sagittis. Integer commodo vitae libero eget sagittis. Etiam id ante dignissim, sollicitudin odio id, rhoncus magna. Vestibulum ante ipsum primis in faucibus orci luctus et ultrices posuere cubilia curae; Praesent at tortor eros. Aenean in euismod mauris. Suspendisse pharetra aliquet ex, sed tempus augue aliquam nec. Donec tempor molestie mi, quis faucibus sem. Cras ornare bibendum nulla a laoreet. Vestibulum risus massa, feugiat at eros ut, rhoncus aliquam turpis.

In malesuada aliquet orci ac rutrum. Donec dignissim lectus eu ante venenatis semper. Aliquam leo ligula, vehicula ac sagittis non, consectetur nec dui. Pellentesque at arcu diam. Aliquam dignissim condimentum auctor. Nulla vehicula et quam vitae scelerisque. Phasellus bibendum posuere egestas. Donec vel nunc massa. Etiam posuere est urna, eu bibendum mauris semper ut. Suspendisse rutrum purus in metus tristique, nec sagittis arcu posuere. Vestibulum tempus pellentesque est sit amet eleifend. Phasellus fringilla neque nec est mattis, sed varius nisi tristique. Pellentesque ullamcorper massa ut orci rutrum pretium.

Nullam tempor elit id lectus maximus, eu congue mauris consequat. Vestibulum facilisis ac lectus eu volutpat. Maecenas ac ante erat. Nunc sit amet velit dapibus, tincidunt risus quis, sagittis est. Aliquam dapibus commodo enim, ut tempor sem accumsan nec. Sed non scelerisque velit. Ut accumsan ullamcorper dui at lacinia.

Morbi blandit molestie lectus non tincidunt. Nam id imperdiet nisi. Proin sed risus vehicula, malesuada tellus nec, dapibus magna. Cras vitae dolor lorem. Morbi vel arcu sit amet dolor auctor facilisis ac nec arcu. Pellentesque congue vulputate ligula, ac tempus metus porttitor sit amet. Aliquam egestas ultricies cursus. Mauris sit amet blandit elit. Cras ac nunc et enim aliquam efficitur. Suspendisse potenti. Integer accumsan turpis ut posuere hendrerit.

Vestibulum vel nisi turpis. Donec pellentesque odio sodales, facilisis elit ac, finibus libero. Maecenas in turpis eget sapien volutpat finibus sed a turpis. In a aliquet magna. Nullam et molestie tellus. In bibendum nunc tortor, eget malesuada turpis tincidunt vitae. Praesent non nulla semper ligula euismod facilisis ac at mauris. Vestibulum ante ipsum primis in faucibus orci luctus et ultrices posuere cubilia curae; Aliquam non faucibus dolor. In hac habitasse platea dictumst. Sed accumsan rutrum fringilla. In ac magna nunc. Praesent ac convallis magna.
}

%=================================================================
%                  Sección editable para siglas
% ================================================================
% ================== NO MODIFICAR LOS COMANDOS ===================
% ================================================================

\itemsiglas{BTU}{Unidad Térmica Británica (\textit{British Thermal Unit})}
\itemsiglas{CRISP-DM}{\textit{Cross-Industry Standard Process for Data Mining}}
\itemsiglas{FIA}{Facultad de Ingeniería y Arquitectura}
\itemsiglas{IA}{Inteligencia Artificial}
\itemsiglas{LLM}{Modelos de Lenguaje de Gran Escala (\textit{Large Language Models})}
\itemsiglas{NTC}{Coeficiente de Temperatura Negativo (\textit{Negative Temperature Coefficient})}
\itemsiglas{OCR}{Reconocimiento Óptico de Caracteres (\textit{Optical Character Recognition})}
\itemsiglas{PLN}{Procesamiento de Lenguaje Natural}
\itemsiglas{RAG}{Generación Aumentada por Recuperación (\textit{Retrieval-Augmented Generation})}
\itemsiglas{UCA}{Universidad Centroamericana José Simeón Cañas}


%=================================================================
%                 Sección editable para abreviaturas
% ================================================================
% ================== NO MODIFICAR LOS COMANDOS ===================
% ================================================================

\itemabreviatura{Abreviaturas}{}

%=================================================================
%               Sección editable para nomenclatura
% ================================================================
% ================== NO MODIFICAR LOS COMANDOS ===================
% ================================================================

\itemnomenclatura{Nomenclatura}{}

%=================================================================
%          Sección editable para capítulo de introducción
% ================================================================
% ================== NO MODIFICAR LOS COMANDOS ==================
% ===============================================================

\newenvironment{planteamientodelproblema}{}{
    En los entornos de soporte técnico y mantenimiento de sistemas de aire acondicionado tipo Split, gran parte del conocimiento especializado necesario para resolver fallas precisas se encuentra distribuido de manera dispersa en extensos manuales de servicio, reportes de fallas, bitácoras y diversa documentación técnica. El acceso manual a esta información para diagnosticar componentes como tarjetas electrónicas, sensores o sistemas de control es un proceso que suele ser lento, inconsistente y dependiente de la experiencia individual de cada técnico.

    De acuerdo con las entrevistas realizadas a profesionales del sector\footnote{Las transcripciones completas de las entrevistas se encuentran en el \hyperlink{anexo:B}{Anexo B}.}, la búsqueda manual de datos técnicos, diagramas o códigos de error toma en promedio entre uno y más de cinco minutos por consulta. Esta falta de información inmediata retrasa significativamente el tiempo de atención de cada servicio y, en consecuencia, impacta de forma directa y negativa en la productividad, reduciendo la cantidad de mantenimientos o reparaciones que los profesionales pueden concretar durante su jornada diaria.

    Ante esta dificultad, los técnicos recurren a fuentes no estructuradas como foros en línea y herramientas de Inteligencia Artificial (IA) genéricas para agilizar sus diagnósticos. Sin embargo, estas alternativas presentan una limitante crítica: las respuestas suelen ser demasiado generales y carecen del contexto técnico específico de cada fabricante. Además, la IA tradicional frecuentemente redunda en la información que entrega, manteniéndose en un círculo de explicaciones superficiales sin profundizar en la solución real requerida para el equipo. Existe, por tanto, un vacío tecnológico claro entre la necesidad de inmediatez del técnico y la exactitud que exigen los manuales de servicio oficiales.
}

\newenvironment{antecedentes}{
Las entrevistas realizadas a cuatro profesionales con perfiles variados y entre seis meses y ocho años de experiencia ---con formación tanto técnica formal como empírica--- validaron la magnitud del problema descrito (véase \hyperlink{anexo:B}{Anexo B}). La totalidad de los entrevistados ha recurrido a herramientas de inteligencia artificial para apoyar sus diagnósticos, calificándolas entre siete y diez sobre diez en confiabilidad. No obstante, a pesar de esta valoración positiva, coinciden en señalar sus limitaciones de forma contundente: un técnico con tres años y medio de experiencia que paga una suscripción a una herramienta de IA y le carga manuales directamente reportó que la herramienta ``se mantiene en el mismo círculo sin profundizar más'', mientras que otro profesional con tres años de trayectoria advirtió que actuar con información descontextualizada respecto al fabricante genera diagnósticos erróneos, evidenciando un riesgo operativo real.

Este déficit también se refleja en la productividad: el técnico con mayor carga de trabajo atiende entre siete y ocho equipos diarios y señaló que los retrasos por búsqueda de información afectan considerablemente los servicios siguientes, en tanto que otro entrevistado indicó que, en casos de alta complejidad, la ausencia de datos técnicos a la mano puede reducir la atención a un único equipo en toda la jornada.

Ante este escenario, los cuatro profesionales manifestaron de forma unánime su interés en un asistente especializado que muestre el fragmento exacto del documento, proporcione pasos numerados e incluya diagramas técnicos; características que, según los entrevistados, eliminarían la ambigüedad de las respuestas genéricas y reducirían el margen de error en diagnósticos que exigen valores precisos de voltaje, resistencia o configuración de sensores. Esta necesidad operativa no resuelta constituye el antecedente directo que motiva el presente proyecto.
}{}

\newenvironment{alcances}{}{
Para dar respuesta a la necesidad operativa de agilizar los diagnósticos técnicos, el proyecto se centra en el diseño y desarrollo de un prototipo funcional basado en la arquitectura de Generación Aumentada por Recuperación (RAG). El ámbito de aplicación abarca exclusivamente los sistemas de aire acondicionado tipo Split de muro o pared, instalados en entornos residenciales y comerciales ligeros. Específicamente, el asistente brindará soporte para equipos con tecnología Inverter y convencional que operen dentro de un rango de capacidad térmica de 9,000 a 24,000 BTU.

A nivel de hardware, el marco de diagnóstico de la herramienta contempla el análisis de los componentes principales que estructuran estos equipos: la unidad interior, la unidad exterior, los sistemas de control, los sensores de temperatura y el sistema de drenaje. Para asegurar la viabilidad del desarrollo dentro de los plazos establecidos, el dominio operativo se acota a la resolución de exactamente diecinueve fallas comunes y repetibles asociadas a dichos componentes, cuyo listado exhaustivo y categorización se documentan en el Anexo A.

El motor de conocimiento del sistema se alimenta mediante un corpus técnico compuesto en su totalidad por manuales de servicio oficiales. Como resultado final, se construye un prototipo de software funcional, accesible a través de una interfaz de usuario (cuyo diseño preliminar se ilustra en el \hyperlink{anexo:C}{Anexo C}), capaz de procesar consultas formuladas en lenguaje natural. Este aplicativo genera hipótesis de causa, secuencias de pasos para el diagnóstico y recomendaciones de mantenimiento básico, garantizando en todo momento la referenciación exacta del fragmento documental de origen que fundamenta la respuesta.
}
\newenvironment{limitaciones}{}{
    \begin{itemize}
        \item \textbf{No entrenamiento de modelos fundacionales:} El sistema no entrenará Modelos de Lenguaje de Gran Escala (LLM) desde cero; la solución integrará modelos preentrenados mediante el uso de frameworks existentes.
        \item \textbf{Ausencia de autonomía de diagnóstico:} El sistema no actuará como una entidad de diagnóstico autónoma; operará estrictamente como un asistente de apoyo para acelerar la búsqueda de información, mientras que la toma de decisiones finales recaerá siempre en el criterio del técnico humano.
        \item \textbf{Respaldo documental obligatorio:} El prototipo no generará recomendaciones empíricas ni respuestas que carezcan de un respaldo documental explícito dentro de la base de datos vectorial.
    \end{itemize}
}
\newenvironment{objetivogeneral}{}{
    Implementar un prototipo funcional del asistente que permita la interacción directa con el técnico.
}
\newenvironment{objetivosespecificos}{}{
    \begin{itemize}
        \item Definir el catálogo de fallas con el que se construirá el sistema RAG.
        \item Diseñar el proceso de diagnóstico de fallas, incluyendo los prompts iniciales y el flujo de conversación esperado.
        \item Analizar el proceso de diagnóstico de fallas en aires acondicionados tipo Split para alinear las respuestas del sistema con el flujo de trabajo real de los técnicos.
        \item Construir un repositorio documental curado a partir de manuales técnicos de servicio, reportes de fallas y registros históricos de los fabricantes.
        \item Integrar un modelo de lenguaje que genere recomendaciones basadas exclusivamente en la información recuperada de las fuentes documentales oficiales.
        \item Evaluar el sistema en términos de utilidad, precisión percibida y reducción del esfuerzo y tiempo de búsqueda manual frente a los métodos tradicionales.
        \item Diseñar el instrumento de evaluación para la encuesta a los usuarios finales.
    \end{itemize}
}


%=================================================================
%                Sección editable para capítulos
% ================================================================
% ================= NO MODIFICAR LOS COMANDOS ====================
% ================================================================

\newcommand{\capitulos}{

\capitulo{Estructura del Trabajo}{

El presente documento se organiza en seis capítulos que siguen una progresión lógica desde la identificación del problema hasta las conclusiones derivadas del desarrollo del prototipo. A continuación se describe brevemente el contenido de cada uno de ellos.

El Capítulo 1, Introducción, expone el planteamiento del problema que motiva el proyecto, describiendo las dificultades que enfrentan los técnicos de mantenimiento al acceder a información especializada de manera oportuna. Se presentan también los antecedentes obtenidos a partir de entrevistas con profesionales del sector, el alcance y las limitaciones operativas del sistema, así como el objetivo general y los objetivos específicos que guían el desarrollo.

El Capítulo 2, Estructura del Trabajo, corresponde al presente capítulo y tiene como propósito orientar al lector sobre la organización del documento, describiendo de forma sintética el contenido de cada sección y la relación entre ellas.

El Capítulo 3, Marco Teórico, fundamenta conceptualmente el proyecto. Se abordan los principios del Procesamiento de Lenguaje Natural (PLN), los Modelos de Lenguaje de Gran Escala (LLM) y su arquitectura Transformer, el problema de la alucinación en sistemas generativos, la representación semántica mediante embeddings, las bases de datos vectoriales y la búsqueda por similitud, la arquitectura de Generación Aumentada por Recuperación (RAG) y el framework LlamaIndex utilizado como capa de integración. Este capítulo incluye una figura explicativa del flujo completo del sistema RAG.

El Capítulo 4, Metodología, describe el proceso de desarrollo del prototipo estructurado bajo la metodología CRISP-DM. Se documenta cada una de sus seis fases: comprensión del problema, comprensión de los datos, preparación de los datos, modelado, evaluación e implementación. Se detallan las decisiones técnicas tomadas en el benchmarking, incluyendo la selección del modelo de embeddings Gemini Embedding 001, el modelo generativo Gemini 1.5 Pro y la base de datos vectorial ChromaDB, así como las estrategias de ingesta y particionado de los manuales de servicio.

El Capítulo 5, Presentación, Análisis e Interpretación de Resultados, expone los hallazgos obtenidos durante la fase de evaluación del prototipo. Se presenta la tabla resumen de pruebas técnicas, las lecciones aprendidas derivadas de errores y hallazgos durante el desarrollo, y los resultados de la evaluación de la herramienta por parte de los técnicos usuarios.

El Capítulo 6, Conclusiones y Recomendaciones, sintetiza los principales logros del proyecto en relación con los objetivos planteados y propone líneas de trabajo futuro sustentadas en bibliografía científica, tales como el uso de bases de datos vectoriales avanzadas, técnicas de expansión de consultas y representación mediante grafos de conocimiento.

Finalmente, el documento incluye dos anexos: el Anexo A contiene el catálogo exhaustivo de las diecinueve fallas documentadas que conforman el dominio del sistema, y el Anexo B presenta el instrumento de evaluación y las transcripciones de las entrevistas realizadas a los técnicos durante la fase de comprensión del problema.

}

\capitulo{Marco Teórico}{

Este capítulo presenta los fundamentos teóricos y las tecnologías subyacentes que sustentan el diseño e implementación del asistente inteligente. Se abordan los conceptos de Procesamiento de Lenguaje Natural, los Modelos de Lenguaje de Gran Escala y su arquitectura Transformer, el problema de la alucinación, la representación semántica mediante embeddings, las bases de datos vectoriales, la arquitectura de Generación Aumentada por Recuperación (RAG) y el framework LlamaIndex utilizado como capa de integración.

\section{Procesamiento de Lenguaje Natural (PLN)}

El Procesamiento de Lenguaje Natural (PLN) es una subdisciplina de la Inteligencia Artificial que estudia la interacción entre las computadoras y el lenguaje humano. Su objetivo principal es permitir que los sistemas computacionales comprendan, interpreten y generen lenguaje natural de forma significativa. El PLN combina técnicas de lingüística computacional, aprendizaje automático, modelado estadístico y redes neuronales profundas \parencite{jurafsky2023speech}.

En el contexto de sistemas técnicos, el PLN permite analizar documentación especializada (manuales de servicio, reportes técnicos, bitácoras) y convertir texto no estructurado en información procesable por un sistema inteligente.

Para el presente proyecto, el PLN es la base que permite interpretar preguntas formuladas por técnicos en lenguaje natural, comprender términos técnicos como códigos de error, voltajes, resistencias o configuraciones de sensores, y relacionar dichas consultas con fragmentos específicos de manuales oficiales.

\section{Modelos de Lenguaje de Gran Escala (LLM) y Arquitectura Transformer}

Los Modelos de Lenguaje de Gran Escala (Large Language Models, LLM) son modelos de aprendizaje profundo entrenados con grandes volúmenes de texto para predecir y generar lenguaje coherente \parencite{brown2020language}.

La mayoría de los LLM modernos están basados en la arquitectura Transformer, introducida por \textcite{vaswani2017attention} en el artículo ``Attention Is All You Need''. Esta arquitectura utiliza un mecanismo llamado autoatención (self-attention) que permite al modelo evaluar la relevancia de cada palabra respecto a las demás dentro de una secuencia.

Entre las características principales de los Transformers se encuentran el procesamiento paralelo de secuencias, la captura eficiente de dependencias de largo alcance y una mejor comprensión contextual que los modelos recurrentes tradicionales.

Sin embargo, los LLM presentan limitaciones importantes: pueden generar información incorrecta (alucinaciones), no garantizan precisión técnica en dominios especializados y no tienen acceso dinámico a documentación específica salvo que se integre un sistema de recuperación externa.

En el contexto del diagnóstico de equipos tipo Split, depender únicamente de un LLM generalista puede producir respuestas ambiguas o no alineadas con especificaciones del fabricante.

\subsection{Limitaciones de los LLM Tradicionales y Problema de Alucinación}

Uno de los principales problemas documentados en los LLM es la alucinación, fenómeno en el que el modelo genera información plausible pero incorrecta o no respaldada por una fuente verificable \parencite{ji2023survey}.

En entornos técnicos críticos, como el diagnóstico eléctrico o electrónico, este comportamiento representa un riesgo operativo real, ya que decisiones incorrectas pueden provocar daños en componentes, generar diagnósticos erróneos, incrementar tiempos de reparación y afectar la seguridad del técnico.

Diversos estudios recientes \parencite{ji2023survey, lewis2020retrieval} han señalado que los LLM sin acceso a fuentes externas estructuradas tienden a completar respuestas con conocimiento estadístico aprendido durante el entrenamiento, lo cual no garantiza exactitud técnica específica por fabricante.

Por ello, surge la necesidad de integrar mecanismos de recuperación de información controlada.

\section{Embeddings y Representación Semántica}

Los embeddings son representaciones vectoriales de palabras, frases o documentos en un espacio matemático de múltiples dimensiones. Su propósito es capturar relaciones semánticas entre conceptos \parencite{mikolov2013efficient}.

En lugar de representar texto como simples palabras clave, los embeddings permiten que términos conceptualmente similares se ubiquen cerca en el espacio vectorial. Por ejemplo, términos como ``Sensor NTC'', ``Termistor'' y ``Sensor de temperatura evaporador'', aunque no sean idénticos, pueden tener proximidad semántica en el espacio vectorial, permitiendo búsquedas más inteligentes.

Esta representación es clave para la búsqueda semántica, que supera la búsqueda tradicional basada en coincidencia exacta de palabras.

\section{Bases de Datos Vectoriales y Búsqueda Semántica}

Las bases de datos tradicionales (relacionales o NoSQL) están diseñadas para búsquedas estructuradas basadas en campos específicos o coincidencias exactas.

En contraste, las bases de datos vectoriales almacenan embeddings y permiten búsquedas por similitud utilizando métricas como la similitud coseno, la distancia euclidiana y el producto punto. Esto permite encontrar fragmentos de documentos que sean semánticamente relevantes aunque no compartan exactamente las mismas palabras.

En el caso de manuales técnicos extensos, una base vectorial permite dividir documentos en fragmentos, convertirlos en embeddings y recuperar únicamente los fragmentos más relevantes según la consulta del técnico. Esta tecnología es esencial para implementar sistemas de recuperación eficiente sobre documentación especializada.

\section{Generación Aumentada por Recuperación (RAG)}

La Generación Aumentada por Recuperación (Retrieval-Augmented Generation, RAG) es una arquitectura que combina un sistema de recuperación de información (base de datos vectorial) con un modelo generativo (LLM) \parencite{lewis2020retrieval}.

En lugar de depender únicamente del conocimiento interno del modelo, RAG permite que el LLM recupere primero fragmentos relevantes de documentos oficiales y genere una respuesta basada exclusivamente en ese contexto recuperado. Esto reduce significativamente el problema de alucinación y aumenta la precisión técnica.

En el proyecto propuesto, el flujo es el siguiente: el técnico formula una consulta, el sistema genera el embedding de la consulta, se buscan fragmentos relevantes en la base vectorial, el LLM genera una respuesta utilizando únicamente esos fragmentos y, finalmente, se muestra el fragmento exacto del manual como respaldo.

Este enfoque permite cumplir con el requisito identificado en las entrevistas: mostrar información específica del fabricante y pasos numerados directamente desde el manual. Como se ilustra en la \reffig{diagrama1.png}, el flujo del sistema se divide en una etapa de recuperación semántica y una etapa de generación fundamentada, ambas orquestadas por el framework LlamaIndex como capa de integración.

\imagen{diagrama1.png}{Arquitectura del Sistema de Generación Aumentada por Recuperación (RAG)}{false}{Elaboración propia}{0.35}{V}

\section{Framework LlamaIndex}

LlamaIndex es un framework diseñado para facilitar la integración entre fuentes de datos externas y modelos LLM \parencite{llamaindex2024}. Su función principal es ingestar documentos, fragmentarlos, generar embeddings, indexarlos y conectarlos con modelos generativos.

De este modo, LlamaIndex permite construir aplicaciones basadas en arquitectura RAG de manera estructurada, facilitando la indexación de manuales técnicos, el control sobre las fuentes utilizadas y la generación de respuestas respaldadas por contexto verificable.

En este proyecto, LlamaIndex actúa como la capa intermedia que conecta los manuales de servicio con el modelo generativo y la base vectorial.

}

\capitulo{Metodología}{

\section{Enfoque Metodológico}

El presente proyecto de graduación tiene un enfoque académico y exploratorio, sin fines de implementación inmediata. No obstante, el sistema se diseña bajo criterios de escalabilidad y adaptabilidad, permitiendo su potencial aplicación en entornos empresariales relacionados con el mantenimiento de equipos de aire acondicionado tipo split.

Dada la naturaleza del presente proyecto, orientado al desarrollo de un sistema basado en generación aumentada por recuperación (RAG), resulta necesario establecer un enfoque metodológico que permita estructurar de manera sistemática las etapas de análisis, diseño y evaluación del sistema.

En este contexto, se selecciona la metodología Cross-Industry Standard Process for Data Mining (CRISP-DM), ampliamente utilizada en proyectos de minería de datos y ciencia de datos \parencite{chapman2000crisp}, y reconocida por su enfoque iterativo y su énfasis en la comprensión del dominio del problema \parencite{shearer2000crisp}.

Si bien no se identifican antecedentes específicos que documenten la aplicación de CRISP-DM en sistemas RAG, sus principios resultan compatibles con este tipo de soluciones, particularmente en lo relativo a la integración de datos, la construcción iterativa de modelos y la evaluación progresiva de resultados \parencite{mariscal2010survey}. Asimismo, los sistemas RAG combinan técnicas de recuperación de información y generación de lenguaje natural para resolver tareas intensivas en conocimiento, lo cual requiere un enfoque metodológico estructurado \parencite{lewis2020retrieval}.

\section{Comprensión del Problema}

Esta fase inicial se orientó a la comprensión del flujo de trabajo real y de las necesidades de los usuarios finales, con el propósito de delimitar el dominio del proyecto. Para ello, se diseñó un instrumento de investigación y se realizaron entrevistas a cuatro profesionales técnicos, cuyos niveles de experiencia oscilaron entre seis meses y ocho años, incluyendo tanto formación técnica formal como empírica.

El análisis de la información recolectada evidenció que la búsqueda manual de información técnica, como códigos de error, diagramas y procedimientos de diagnóstico, en documentación física o digital, es la etapa clave para la realización de un servicio de mantenimiento \parencite{krutli2025augmenting}, estimando que se requiere entre uno y más de cinco minutos por consulta puntual. Este tiempo impacta directamente en la eficiencia operativa, generando retrasos en la prestación del servicio y una disminución en la productividad diaria.

Asimismo, se identificó que los técnicos recurren a herramientas de inteligencia artificial de propósito general, como ChatGPT, con el objetivo de agilizar sus procesos de diagnóstico. No obstante, dichas herramientas presentan limitaciones significativas \parencite{anhhoang2025survey}, tales como la generación de respuestas redundantes, ambiguas o carentes de contexto técnico específico por fabricante, lo cual introduce un riesgo operativo en escenarios reales de mantenimiento .

A partir de estos hallazgos, se delimitó el dominio del proyecto al diagnóstico preliminar y mantenimiento de sistemas de aire acondicionado tipo split de muro, enfocados en entornos residenciales y comerciales ligeros, con capacidades entre 9,000 y 24,000 BTU, abarcando tanto tecnologías Inverter como convencionales. El alcance funcional se restringió a un conjunto estructurado de diecinueve fallas comunes y recurrentes, asociadas a los principales componentes del sistema.

\section{Comprensión de los Datos}

En esta fase se llevó a cabo la recolección, evaluación y clasificación de la base de conocimiento destinada a alimentar el sistema RAG. La estrategia de búsqueda se enfocó exclusivamente en manuales de servicio oficiales y reportes técnicos de fabricantes reconocidos, descartando los manuales de usuario convencionales debido a su limitada profundidad en procesos de diagnóstico.

Con el objetivo de garantizar la calidad y relevancia de la base de conocimiento documental, se definió un protocolo de selección que exigía la presencia de secciones explícitas de resolución de problemas (troubleshooting), así como diagramas de flujo que establecieran relaciones claras entre síntomas, causas y acciones correctivas.

Adicionalmente, se excluyó documentación correspondiente a sistemas industriales, chillers y equipos portátiles, con el fin de evitar la introducción de información no pertinente que pudiera afectar la precisión del sistema.

Como criterio técnico adicional, se estableció la necesidad de legibilidad computacional de los documentos, requiriendo que estos se encontrarán en formato PDF nativo o que incluyeran capas de texto seleccionable. En consecuencia, se descartaron documentos complejos escaneados sin procesamiento de reconocimiento óptico de carácteres (OCR), debido a su incompatibilidad con los procesos de generación de representaciones vectoriales (embeddings).

Finalmente, los documentos seleccionados fueron sometidos a un esquema de evaluación cuantitativa en una escala de 0 a 100, priorizando aquellos que incluían códigos de error específicos, diagramas eléctricos detallados y tablas técnicas relevantes, considerados como elementos críticos para el proceso de diagnóstico.

\section{Modelado de Datos}

En esta fase se llevó a cabo el diseño, construcción y configuración del motor de la solución tecnológica, integrando los componentes de la arquitectura de RAG mediante el uso del framework LlamaIndex como orquestador principal.

A diferencia de un enfoque estático, la configuración del sistema fue determinada a partir de un proceso iterativo de experimentación y validación empírica, sustentado en la ejecución de un plan de pruebas técnicas sobre el pipeline completo. Dichas pruebas abarcaron desde la ingesta de documentos hasta la generación final de respuestas, permitiendo identificar la interdependencia entre las decisiones de diseño en cada etapa del sistema .

Los resultados evidenciaron que las decisiones tomadas en fases previas, particularmente en la preparación de los datos, tienen efectos en cascada sobre la calidad de los embeddings y la precisión de la recuperación, validando así el enfoque iterativo propuesto por la metodología CRISP-DM .

En el componente de representación vectorial, se llevó a cabo un proceso de benchmarking comparando modelos de código abierto y soluciones propietarias. Los hallazgos demostraron que los modelos monolingües en inglés presentan un desempeño limitado en escenarios bilingües, debido a su incapacidad para mapear consultas en español con documentación técnica en inglés. En este contexto, el modelo gemini-embedding-001 fue seleccionado por su capacidad multilingüe, alcanzando un desempeño significativamente superior en términos de recuperación semántica \parencite{lee2025gemini}.

La información vectorizada fue almacenada utilizando ChromaDB como base de datos vectorial de código abierto caracterizada por un buen balance entre velocidad de recuperación, calidad de embeddings y rendimiento (Ammar, 2025), permitiendo la indexación eficiente de los documentos procesados.

En cuanto al componente de recuperación, se evaluaron diferentes configuraciones de búsqueda y síntesis de respuesta. Los resultados mostraron que el uso de estrategias como tree\_summarize proporciona una mayor coherencia lógica en consultas de diagnóstico complejo, mientras que el modo compact resulta más eficiente para consultas directas relacionadas con códigos de error .

Finalmente, se integró el modelo generativo Gemini 1.5 Pro para la construcción de las respuestas \parencite{reid2024gemini}, configurado con parámetros orientados a la precisión y reducción de variabilidad. La configuración final del sistema representa el resultado de un proceso de optimización basado en evidencia experimental, en el cual cada componente fue seleccionado en función de su impacto medible en la calidad de las respuestas generadas.

De forma paralela a la configuración del motor lógico del sistema RAG, se desarrolló el diseño inicial de la interfaz gráfica web destinada a la interacción con el usuario final. Este prototipo contempla dos componentes funcionales principales: un sistema de chat conversacional, mediante el cual el técnico formula consultas en lenguaje natural y recibe respuestas fundamentadas en la base vectorial, y un módulo de gestión documental que permite la carga y administración de los manuales de servicio que alimentan el corpus técnico del sistema. Las capturas de pantalla correspondientes al diseño preliminar de esta interfaz se documentan como evidencia en el \hyperlink{anexo:C}{Anexo C}.

}

\capitulo{Presentación, Análisis e Interpretación de Resultados}{

}
}

%=================================================================
%       Sección editable para conclusiones y recomendaciones
% ================================================================
% ================= NO MODIFICAR LOS COMANDOS ===================
% ================================================================

\newenvironment{conclusiones}{
    
}{}

\newenvironment{recomendaciones}{
       
}{}

%=================================================================
%                Sección editable para glosario
% ================================================================
% ================= NO MODIFICAR LOS COMANDOS ====================
% ================================================================

\newcommand{\mostrarglosario}{false}

%=================================================================
%                    Sección editable para anexos
% ===============================================================
% ================== NO MODIFICAR LOS COMANDOS ==================
% ================================================================

\newcommand{\anexos}{

    \anexo{Catálogo de Fallas Comunes en Aires Acondicionados tipo Split}{
        \tabla{V}{
          \begin{longtable}{|p{7.5cm}|p{7.5cm}|}
            \hline
            \textbf{Síntoma Observable} & \textbf{Posible Causa} \\
            \hline
            El equipo no arranca y muestra error interno & Falla de memoria (EEPROM) en la tarjeta electrónica interior \\
            \hline
            El compresor no inicia y aparece error en unidad exterior & Falla de memoria EEPROM en la tarjeta electrónica exterior \\
            \hline
            El equipo se apaga solo y muestra código de error & Posible fuga de refrigerante o sensor defectuoso \\
            \hline
            El equipo no enfría y muestra error & Sensor de temperatura abierto o en corto \\
            \hline
            Funcionamiento irregular o apagado inesperado & Sensor de temperatura defectuoso \\
            \hline
            Error de lectura de temperatura & Sensor de temperatura dañado \\
            \hline
            El equipo se protege y se detiene & Sensor de temperatura defectuoso \\
            \hline
            Apagado por protección térmica & Sensor de temperatura fuera de rango \\
            \hline
            El equipo se apaga por protección eléctrica & Sobrecorriente o bloqueo del sistema \\
            \hline
            El equipo se detiene inmediatamente después de encender & Protección del módulo de potencia (IPM) \\
            \hline
            El equipo no detecta temperatura ambiente & Sensor ambiente abierto o en corto \\
            \hline
            El equipo no enfría correctamente & Sensor del evaporador defectuoso \\
            \hline
            Ventilador exterior no gira correctamente & Falla del motor del ventilador exterior \\
            \hline
            El display no responde o muestra error & Falla de comunicación entre tarjeta electrónica y display \\
            \hline
            Bajo rendimiento y apagado del sistema & Detección de fuga de refrigerante \\
            \hline
            El equipo no arranca & Falla en módulo de potencia (IPM) \\
            \hline
            El equipo se apaga por protección & Bajo voltaje de alimentación \\
            \hline
            Protección por voltaje alto & Sobrevoltaje en la línea eléctrica \\
            \hline
            El compresor no inicia & Falla de arranque del compresor \\
            \hline
          \end{longtable}
        }
    }

    \anexo{Instrumento de Evaluación y Transcripciones de Entrevistas}{

    \sectionanexo{Entrevista 1: Edwin Roberto Quintanilla}

    \textbf{Nombre:} Edwin Roberto Quintanilla\\

    \textbf{¿Nos da su consentimiento para utilizar sus respuestas con fines académicos en esta investigación?}\\
    Sí, doy mi consentimiento.

    \textbf{¿Cuántos años tiene trabajando en mantenimiento de aires acondicionados tipo Split?}\\
    En refrigeración y mantenimientos de aires llevo 3 años.

    \textbf{¿Cuenta con formación técnica formal o aprendió de manera empírica?}\\
    Los conocimientos que adquirí sobre refrigeración y aires acondicionados tipo mini Split fueron de manera empírica. Ya que en mi anterior trabajo dieron unas capacitaciones con respecto a esta rama de refrigeración, sin embargo, pues no solo me quedé con lo aprendido en esas capacitaciones, sino que a la hora de estar ejecutando pues me tocó ir aprendiendo. Entonces, en esas capacitaciones que fueron informales, se podría decir, no me otorgaron algún título; entonces, por esa manera, sí puedo decir en mi respuesta que ha sido de manera empírica.

    \textbf{¿Trabaja principalmente solo o en equipo?}\\
    Actualmente, sí, trabajo principalmente solo.

    \textbf{¿Atiende múltiples marcas o se especializa en algunas específicas?}\\
    Con respecto a la hora de la manipulación de los aires acondicionados tipo mini Split, ya sea convencional o Inverter, la marca, pues no, no tengo una especialización con una marca en específica. Sin embargo, pues tengo marcas con las que estoy más familiarizado, ¿verdad? Como Midea, York, Andina, LG, Mitsubishi, con las que, pues, en ciertas ocasiones he logrado tener la interacción con ese tipo de equipo, ya sea para un mantenimiento, una instalación o una reparación.

    \textbf{Cuando llega a revisar un equipo que falla, ¿cuál es su secuencia típica de diagnóstico?}\\
    Pues para dar un buen diagnóstico de un aire acondicionado, lo primero que hago es conectar la pinza amperimétrica para ver el consumo del equipo con respecto a lo de la placa del fabricante, ¿verdad? Ya que, pues, cada fabricante tiene sus diferentes especificaciones. Entonces, conectar también los manómetros como segunda parte, para ver en cuántas presiones está el refrigerante, en cuántos PSI, y también comparar los PSI con respecto a lo que dice la placa del equipo, ¿verdad? Ya que cada equipo, pues, dependiendo su capacidad de BTUs, varía la cantidad de refrigerante que tiene que poseer en el sistema, y también depende de la marca, ¿verdad? Ya que otros usan diferentes tipos de refrigerante. Entonces, como segundo paso es eso, ¿verdad? Encontrar, ver las mediciones. Ya después, con esos dos datos, parte el diagnóstico de si un componente eléctrico está dañado. Eso ya uno lo va viendo con respecto a cómo va destapando las unidades, ¿verdad? Tanto la evaporadora como la condensadora.

    Y para seguir agregando más puntos, todo varía con respecto a lo que el cliente da como diagnóstico inicial, ¿verdad? Que el cliente dice: ``el equipo huele mal'', ``el equipo está botando agua'', ``el equipo no está enfriando''. Porque los anteriores pasos que he dicho es cuando el equipo no está enfriando o, pues, ha tenido fallas o no enciende. Ya si está goteando, pues se entiende que cuando un equipo está botando agua es porque no ha tenido su mantenimiento, ya que se recomienda que sea cada seis meses. Ya que, aunque el equipo no se utilice, está la unidad que está en el exterior, entonces pues esa unidad siempre recibe la contaminación del aire, ¿verdad? Uno entiende que esa unidad está a la intemperie, entonces necesita que tenga sus mantenimientos al día para que pueda hacer una buena transmisión de calor.

    Entonces, pues, con respecto a hacer un diagnóstico, lo principal siempre primero es tomar mediciones, ya sea de refrigerante, que sería con los manómetros, y con el amperímetro para ver cuánto está consumiendo y cuánto está llegando de voltaje a los bornes. Ya de ahí se parte a los siguientes desglosamientos de diagnóstico, ¿verdad? Ya que, bueno, algunos diagnósticos en el momento los va a ver uno; qué es lo que se necesita. Bueno, siempre es el que uno logra verificar qué es lo que está pasando, ¿verdad? Ya sea que al ventilador se le quemara su capacitor, o sus bobinas se sobrecalentaran, o que el capacitor pues ya estuviera dañado. O la placa, porque bueno, hoy en día los equipos Inverter pues son demasiado ahorrativos; sin embargo, pues, el hecho de tener tantos sensores y que sean muy electrónicos, al no tener sus protecciones a la hora de tener una buena instalación, pues muchas veces se terminan dañando las placas. Entonces, también a veces hay que inspeccionar los bornes de las placas para ver si también llegan a su voltaje, o si no hay algún componente que se pueda ver a la vista que esté dañado, pues uno con el tester hace las correctas mediciones correspondientes al tipo de componente que uno va a estar verificando.

    \textbf{¿En qué momento del proceso necesita consultar información externa (manual, internet, compañero de trabajo)?}\\
    Con respecto a su pregunta, tengo dos respuestas. Bueno, mi respuesta es para dos opciones, que sería tanto la ayuda de un compañero como la ayuda de Internet. Ya que, pues, el Internet es esencial en este rubro para tener tablas con respecto a cada equipo, qué componentes lleva, el diagrama eléctrico... Ya que muchas veces, si uno no es la persona que ha instalado el aire, pues la gente termina botando los manuales, ¿verdad? Entonces, para adquirir esa información de los manuales sí necesita Internet uno; ya que, pues, si bien el equipo trae sus placas donde tiene esas especificaciones de refrigerante, de voltaje, de consumo, siempre tener el manual de servicio es algo importante.

    Y también menciono que más allá de tener el Internet, que es un aliado demasiado bueno, también pues si uno tiene algún compañero... pues pongo este contexto de que a veces la experiencia de otra persona que ya tuviera fallas parecidas, o pues el conocimiento de años en este rubro, también siento que tiene su ayuda. Que también eso es algo que, pues sí, si muy bien con el Internet uno puede llevarse, pues a veces el hecho de tener a una persona con experiencia en el rubro también va de la mano. Pero también esas personas que ya tienen años en ese rubro no pueden desenvolverse, ¿verdad?, de una manera óptima, se puede decir. Porque, bueno, hoy la tecnología va avanzando demasiado, que pues los refrigerantes cambian, las tecnologías de la electrónica de los aires cambian bastante porque siempre van mejorando. Entonces, si bien es bueno tener ahí a una persona con la cual te puedas apoyar con una pregunta, que te pueda ayudar, pues siempre se necesita el Internet. Ya que, pues, a veces es muy complicado el hecho de tener capacitaciones constantes de los proveedores de los equipos, ya que no siempre las hay, ¿verdad? Entonces, la manera en la cual uno se puede defender bastante con conocimientos que sean correctos (porque en este rubro no se puede equivocar uno), pues es teniendo el Internet de la mano y saber buscar bien los manuales, porque eso es lo que lleva más tiempo.

    \textbf{¿Qué tipo de fallas le obligan a buscar documentación? (Códigos de error, tarjetas electrónicas, sensores, fallas eléctricas, etc.)}\\
    Bueno, de las fallas, se puede decir que más que todo para los códigos de error sí se necesita el Internet. Ya que, por mi caso, al tener un catálogo de varias marcas, pues no todas las veces el error es igual en todas las marcas. Entonces, sí, si bien en algunos errores cambian las palabras o son similares, no siempre el error ``3'' va a ser en Andina y luego va a ser en LG. Es muy difícil a veces que todos los errores sean muy iguales; siempre es mejor tener la base técnica en la cual uno no lo pueda estar, pues, inventando.

    \textbf{Cuando no encuentra la información rápidamente, ¿qué suele ocurrir?}\\
    Lo principal es que el servicio se retrasa. Ya que eso de la prueba y error retrasa todavía aún más, y pues el hecho de interpretar mal un manual... más que todo sería encontrar la información en el manual, porque el hecho de que uno esté interpretando un manual requiere tiempo. Entonces, la verdad es que el hecho de no tener el conocimiento a la mano, pues sí hace que un trabajo se extienda demasiado tiempo y pues ya no sea un trabajo que se puede ejecutar rápido, sino que por lo menos una hora, por decir un tiempo. Por el hecho de que, pues, siempre hay que tener la información correcta, no se puede buscar o usar prueba y error. Entonces, sí se necesita, sí o sí, una fuente muy confiable.

    \textbf{¿Le ha ocurrido interpretar mal un manual?}\\
    En mi caso no me ha pasado por el momento.

    \textbf{¿Qué tan fáciles de entender considera que son los manuales de servicio?}\\
    Siempre, pues, el manual es cuestión de perspectiva, ¿verdad? Ya que hay manuales que son de diferentes marcas, más ---se puede decir--- accesibles a la interpretación, ¿verdad? Entonces siento que, en ese caso, ya sea a la hora de leer un manual de servicio, siempre más allá de que sea fácil o difícil entenderlo, requiere su tiempo. Ya que, pues, estar buscando justamente en el manual que es para el equipo la pregunta, consulta o el dato técnico que se necesita, pues sí requiere un esfuerzo, ¿verdad? Un esfuerzo y tiempo para poder tener una conclusión buena.

    \textbf{Cuando necesita información técnica, ¿qué consulta primero?}\\
    El principal sería Google, ya que pues allí está la mayoría, o si no es que todos los manuales de servicio. Pero es bastante lento, ya que el hecho de encontrar la marca, el manual, la página y el apartado de lo que estás buscando es un proceso que sí necesita su tiempo. Y que también viene quedando bastante lento, como YouTube; ya que, pues sí hay muchos videos empíricos que explican diferentes cuestiones con respecto a la refrigeración. Volvemos a que no es una base técnica como uno esperaría en un manual de servicio. Entonces sí, los dos son pues bastante eficientes en su manera... bueno, no se puede decir eficiente porque lleva bastante tiempo, pero los dos ayudan. Sin embargo, no dejan siempre eso que se desea para que uno pueda facilitarse más.

    \textbf{En promedio, ¿cuánto tiempo le toma encontrar la información cuando debe buscarla?}\\
    Definitivamente sí, serían más de cinco minutos.

    \textbf{Cuando necesita buscar información técnica y eso le toma varios minutos, ¿ese tiempo impacta la cantidad de servicios que puede realizar en el día?}\\
    La falta de información inmediata definitivamente perjudica el tiempo de la ejecución de los trabajos. Ya que así como se pueden ver en un día tres equipos, solamente se puede llegar a ver uno por el hecho de que muchas veces el tener un dato técnico del equipo, de la ficha técnica, o que se pueda llegar a encontrar el manual de servicio... y no tener eso a la mano, puede retrasar toda la operación.

    \textbf{Aproximadamente, ¿cuántos equipos atiende en promedio por día?}\\
    El entrevistado no proporcionó una respuesta específica.

    \textbf{¿Ha utilizado herramientas como ChatGPT u otras Inteligencias Artificiales para resolver fallas técnicas?}\\
    Sí, actualmente pues he intentado implementar poco a poco la inteligencia artificial para poder ayudarme en el día a día con el trabajo.

    \textbf{En una escala del 1 al 10, ¿qué tan confiables han sido las respuestas de la IA?}\\
    Sería un 10, ya que la inteligencia artificial es una herramienta que hoy en día está en una constante evolución y a pasos agigantados, y el hecho de poderla implementar en mi trabajo ha sido algo estupendo. Ya que, al poder brindarle el contexto y también saber cómo preguntar o qué es lo que estoy preguntando, podemos entender cómo llegar a una respuesta con esa confianza total. Entonces sí, está muy bien hoy en día, pues, tener que la inteligencia artificial te ayude en algún tema. Bueno, en mi caso, en mi trabajo, me resulta a veces muy difícil, ya que pues me toma tiempo también el hecho de poder utilizarla de momento, porque tengo que... bueno, sería bueno que ella estuviera de la mano con el rubro de mi trabajo. Entonces, me ayuda bastante y puedo darle la confianza de 10; sin embargo, pues todavía siento que tuviera que trabajar más para llegar a la evolución que yo quisiera para poder implementarla en mi trabajo, pero esperaría que en su momento pueda tener algo parecido.

    \textbf{¿Le redujo el tiempo comparado con buscar en un manual?}\\
    Definitivamente sí; en comparación con los manuales, la inteligencia artificial me ha ayudado a ahorrar mucho más tiempo. Si podría decir que andar buscando en un manual podría ser de más de cinco minutos, la inteligencia artificial me acorta el tiempo bastante, bastante, bastante más.

    \textbf{¿Alguna vez le dio una respuesta incorrecta?}\\
    Respuestas incorrectas, mejor dicho, o respuestas con un contexto lejano a la pregunta que uno realiza. Ya que en ese caso siento que al no brindarle todo el contexto de que es un tema específico, como mi trabajo de refrigeración, pues muchas veces se puede quedar un poco corta la idea del conocimiento que brinda o la respuesta que uno necesita, ya que se sale un poco de contexto. Entonces, al tomarse uno el tiempo de poder elaborar unas preguntas y brindar el contexto, todo se soluciona. Sin embargo, lleva su tiempo, pero es mucho más sencillo que llevar una búsqueda en un manual de servicio.

    \textbf{¿Confía más en la IA o en manuales oficiales?}\\
    Pues al principio sí reconozco que siempre consultaba manuales de servicio. Sin embargo, pues, el tiempo de estar utilizando la inteligencia artificial me ha dado la confianza que también me transmiten los manuales de servicio para poder, pues, ejecutar mi trabajo con la inteligencia artificial de la mano. Y, definitivamente, buscar en manuales es mucho más tardado.

    \textbf{¿Confía en la información que encuentra en grupos de técnicos o YouTube?}\\
    Pues la verdad es que en grupos, foros o videos hay mucha información en la que uno, pues, puede confiar porque son de colegas que se dedican al mismo rubro. Sin embargo, pues uno no solo se puede quedar con el hecho de un video, sino que siempre hay que buscar en los manuales de servicio que son donde hay más información, y, en este caso que estaba implementando, usar Gemini para poder tener una respuesta más rápida.

    \textbf{¿Qué le preocupa más del uso de Inteligencia Artificial en su trabajo?}\\
    Lo principal es que es muy general. Ya que uno sí necesita estar enfocado en cada marca en su momento. Entonces la IA, al principio, es muy general y es donde uno puede cometer el error de no darle contexto y llegar a entender algo que no se puede hacer. Y pues si uno en este rubro hace algo que no va con respecto al fabricante que está en el manual de servicio, pues ahí uno termina demasiado mal, en problemas, arruinando el equipo. Entonces, sí, uno puede utilizar la IA, la idea es buena porque es rápida y tiene mucha información, pero también uno tiene que saber usar la herramienta. Y pues también, que la herramienta fuera más especializada en lo que necesito, ¿verdad? En cada marca.

    \textbf{Si tuviera un asistente en su celular que solo responda basándose en Manuales de Servicio oficiales, muestre el fragmento exacto del manual, indique pasos numerados e incluya diagramas eléctricos, ¿le sería útil?}\\
    Claramente sería útil. Ya que en mi rubro, ¿verdad?, que tengo un catálogo de diferentes marcas, el hecho de tener el manual de servicio de cada equipo, de cada marca, no es tan viable... bueno, para nada. Entonces, tener a la palma de mis manos el manual de servicio, o una herramienta que pueda ayudarme a interpretar el manual, o buscarme más allá de interpretar, y luego asegurarse que en el manual está esa información, sería estupendo.

    \textbf{¿Lo usaría durante el servicio en campo o solo para capacitación?}\\
    Es que, el hecho de tener a un asistente tan accesible y con unas funciones demasiado buenas, se tiene que usar sí o sí en el campo y para capacitaciones. Ya que es una herramienta... podría decir que se convertiría en una de las herramientas principales. El hecho de tener ese asistente, como en este rubro, ¿verdad?, es bueno tener un teléfono solo para el trabajo. El hecho es que, tener siempre cargado ese teléfono, es como dejar los manómetros en refrigeración, no se puede olvidar. Entonces, si llega a existir una herramienta con todas esas funciones, sí o sí se tiene que utilizar día a día.

    Y sería una herramienta que a los colegas que están en este rubro de refrigeración les fascinaría, ya que simplemente con el hecho de saber que pudiera haber una inteligencia artificial con esas funciones, es muy increíble.

    \vspace{1cm}

    \sectionanexo{Entrevista 2: Daniel}

    \textbf{¿Cuántos años tiene trabajando en mantenimiento de aires acondicionados tipo Split?}\\
    Aproximadamente 3 años y medio.

    \textbf{¿Cuenta con formación técnica formal o aprendió de manera empírica?}\\
    Formación técnica.

    \textbf{¿Trabaja principalmente solo o en equipo?}\\
    Por lo general solo, aunque a veces se hacen instalaciones de aires en equipo.

    \textbf{¿Atiende múltiples marcas o se especializa en algunas específicas?}\\
    Atiende cualquier tipo de marcas de aires acondicionado.

    \textbf{Cuando llega a revisar un equipo que falla, ¿cuál es su secuencia típica de diagnóstico?}\\
    Menciona que todo depende de la falla que se esté presentando en el equipo, pero por lo general inicia con pruebas de encendido y apagado, pruebas de enfriamiento y procede con pruebas eléctricas.

    \textbf{¿En qué momento del proceso necesita consultar información externa (manual, internet, compañero de trabajo)?}\\
    Por lo general cuando es un error poco usual.

    \textbf{¿Qué tipo de fallas le obligan a buscar documentación? (Códigos de error, tarjetas electrónicas, sensores, fallas eléctricas, etc.)}\\
    Falla de sensores, fallas eléctricas, fallo de algún componente.

    \textbf{Cuando no encuentra la información rápidamente, ¿qué suele ocurrir?}\\
    En general se retrasa el servicio, y dependiendo si no es un error muy complejo procede a hacer prueba y error, siempre manteniendo el cuidado. Mencionó que sí le ha pasado que interpretó mal un manual, ya que al ser algo extensos y confusos, suele pasar que no se interpreta a la perfección.

    \textbf{¿Qué tan fáciles de entender considera que son los manuales de servicio?}\\
    Poco intuitivos y bastante tediosos de leer por lo general.

    \textbf{Cuando necesita información técnica, ¿qué consulta primero?}\\
    Directamente a una IA, en este caso mencionó ChatGPT, y esta persona le pasa los manuales técnicos que usa la empresa.

    \textbf{En promedio, ¿cuánto tiempo le toma encontrar la información cuando debe buscarla?}\\
    En promedio 1--5 minutos usando IA; si es de forma manual, más de 5 minutos.

    \textbf{Cuando necesita buscar información técnica y eso le toma varios minutos, ¿ese tiempo impacta la cantidad de servicios que puede realizar en el día?}\\
    Impacta de forma considerable, ya que retrasa mucho los siguientes servicios. Por lo general hace entre 7--8 servicios por día.

    \textbf{¿Ha utilizado herramientas como ChatGPT u otras IA para resolver fallas técnicas?}\\
    Sí, ChatGPT.

    \textbf{En una escala del 1 al 10, ¿qué tan confiables han sido las respuestas?}\\
    Lo considera un 8 sobre 10, a nivel de confianza.

    \textbf{¿Le redujo el tiempo comparado con buscar en un manual?}\\
    Sin duda, le ahorra muchísimo tiempo.

    \textbf{¿Alguna vez le dio una respuesta incorrecta?}\\
    Completamente incorrecta no, pero a veces es redundante con las respuestas, y se mantiene en el mismo círculo sin profundizar más.

    \textbf{¿Confía más en IA o en manuales oficiales?}\\
    En este caso menciona que en la IA, por el hecho que le pasa directamente el manual técnico y se basa puramente en ello.

    \textbf{¿Confía en la información que encuentra en grupos o YouTube?}\\
    Sí, por lo general al tener el conocimiento, según lo que ha buscado considera que ha sido información correcta y útil.

    \textbf{¿Qué le preocupa más del uso de IA?}\\
    Que brinde información errónea en algún punto.

    \textbf{Si tuviera un asistente en su celular que solo responda basándose en Manuales de Servicio oficiales, muestre el fragmento exacto del manual, indique pasos numerados e incluya diagramas eléctricos, ¿le sería útil?}\\
    Sin duda, menciona que sería muy útil.

    \textbf{¿Lo usaría durante el servicio en campo o solo para capacitación?}\\
    En ambos casos, principalmente más en servicio, ya que simplifica y ahorra mucho tiempo.

    \textbf{¿Estaría dispuesto a pagar por una herramienta así si realmente le ahorra tiempo?}\\
    Sí, el entrevistado menciona que actualmente paga ChatGPT, y sin duda considera que le ahorra bastante tiempo en su trabajo, para ser más productivo, y sin duda pagaría algo que estuviera mucho más especializado en dicho campo.

    \vspace{1cm}

    \sectionanexo{Entrevista 3: Técnico con formación en ingeniería mecánica}

    \textbf{¿Cuántos años tiene trabajando en mantenimiento de aires acondicionados tipo Split?}\\
    He trabajado en el rubro por 6 meses.

    \textbf{¿Cuenta con formación técnica formal o aprendió de manera empírica?}\\
    Formación universitaria en ingeniería mecánica.

    \textbf{¿Trabaja principalmente solo o en equipo?}\\
    Usualmente en equipo.

    \textbf{¿Atiende múltiples marcas o se especializa en algunas específicas?}\\
    He trabajado múltiples marcas y diferentes dispositivos.

    \textbf{Cuando llega a revisar un equipo que falla, ¿cuál es su secuencia típica de diagnóstico?}\\
    Verificar la falla reportada, realizar preguntas, consultar información o manuales y solucionar el problema.

    \textbf{¿En qué momento del proceso necesita consultar información externa (manual, internet, compañero de trabajo)?}\\
    Cuando ya no encuentro la información por mi cuenta o necesito una segunda opinión.

    \textbf{¿Qué tipo de fallas le obligan a buscar documentación? Cuando no encuentra la información rápidamente, ¿qué suele ocurrir?}\\
    Cuando los dispositivos muestran códigos de error me veo forzado a consultar en internet. El servicio puede retrasarse dependiendo de la complejidad y a veces toca hacer prueba y error, los manuales se interpretan bien sin tanta complicación.

    \textbf{¿Qué tan fáciles de entender considera que son los manuales de servicio?}\\
    Son fáciles de entender, pero muchos manuales considero que están incompletos o sería bueno que el fabricante diera más información.

    \textbf{Cuando necesita información técnica, ¿qué consulta primero?}\\
    Si tuviera que ponerles un orden: primero YouTube, segundo grupos de técnicos, luego manuales de servicios y Google.

    \textbf{En promedio, ¿cuánto tiempo le toma encontrar la información cuando debe buscarla?}\\
    Por lo general suele durar de 1 a 5 minutos.

    \textbf{Cuando necesita buscar información técnica y eso le toma varios minutos, ¿ese tiempo impacta la cantidad de servicios que puede realizar en el día? Aproximadamente, ¿cuántos equipos atiende en promedio por día?}\\
    No impacta mucho, a menos que verdaderamente sea algo muy complejo. Por lo general si es mantenimiento correctivo, 3 equipos y si son rutinas preventivas o correctivas leves ya programadas pueden ser 6 equipos.

    \textbf{¿Ha utilizado herramientas como ChatGPT u otras IA para resolver fallas técnicas?}\\
    Sí las he utilizado para resolver fallas técnicas. Por lo general, las respuestas son confiables; en una escala de 1 al 10 diría que 7. Sí ayuda, es menor el tiempo a manualmente consultar un documento o manual, pero siempre hay que validar la información.

    \textbf{¿Confía en la información que encuentra en grupos o YouTube? ¿Qué le preocupa más del uso de IA?}\\
    Sí confío en la información de grupos y YouTube. Me preocupa que sea una respuesta muy general y no abunde en la información.

    \textbf{Si tuviera un asistente en su celular que solo responda basándose en Manuales de Servicio oficiales, muestre el fragmento exacto del manual, indique pasos numerados e incluya diagramas eléctricos, ¿le sería útil?}\\
    Que el asistente indique pasos numerados de cómo solucionar o dar mantenimiento, e incluya diagramas eléctricos.

    \textbf{¿Lo usaría durante el servicio en campo o solo para capacitación?}\\
    Lo usaría para ambos, tal vez más para capacitación.

    \textbf{¿Estaría dispuesto a pagar por una herramienta así si realmente le ahorra tiempo?}\\
    No, creo que es un poco exagerado ya que la verdad no realiza mucho para que sea justificado pagar por ello.

    }

    \anexo{Diseño Preliminar de la Interfaz de Usuario}{

    A continuación se presentan las capturas de pantalla correspondientes al diseño preliminar de la interfaz gráfica web del sistema de Generación Aumentada por Recuperación. Estas evidencias documentan los dos componentes funcionales principales del prototipo: el sistema de chat conversacional para consultas de diagnóstico y el módulo de gestión documental para la carga de manuales técnicos.

    \imagenanexo{frontend1.jpeg}{Interfaz principal del asistente de diagnóstico}{0.14}{V}

    \imagenanexo{frontend2.jpeg}{Módulo de gestión y carga de manuales técnicos}{0.14}{V}

    }

}